\documentclass[11pt,a4paper]{article}

\usepackage{amsmath,amssymb}
\usepackage{graphicx}
\usepackage{geometry}
\usepackage{hyperref}
\usepackage{booktabs}
\geometry{margin=1in}

\title{BIMOD: A Nasion-Based Polarity Discrimination Model Toward Polarity-Aware Instrumentation}
\author{Author Name}
\date{}

\begin{document}

\maketitle

\begin{abstract}
This study proposes BIMOD, a nasion-centered structural polarity discrimination model formalized through vector-field abstraction and energy-based resonance conditions. The framework defines an eight-position structural polarity mapping and introduces a testable constraint formulation under information-mediated conditions. An auxiliary double-blind optical protocol is outlined to evaluate classification consistency using raw, unprocessed photographic data. The model is presented as a formal abstraction requiring empirical validation and is positioned as a step toward polarity-aware measurement architectures.
\end{abstract}

\section{Introduction}

This study is motivated by the objective of advancing human-centered measurement, understood as measurement architectures that formally incorporate structured individuality. 

The present work proposes a nasion-based polarity discrimination model (BIMOD) framed as a formal and testable abstraction. No new physical law is claimed; empirical validation remains necessary.

\section{Structural Polarity Mapping}

An eight-position structural polarity mapping is defined relative to the nasion reference point. Let the polarity state be represented as:

\[
P_i \in \{+1, -1\}, \quad i = 1,\dots,8
\]

The structural configuration is represented as a vector field:

\[
\mathbf{F} = \sum_{i=1}^{8} P_i \mathbf{e}_i
\]

where $\mathbf{e}_i$ are basis vectors corresponding to defined spatial orientations.

\section{Energy-Based Resonance Condition}

We define an interaction energy functional:

\[
E = - \sum_{i=1}^{8} P_i S_i
\]

where $S_i$ denotes stimulus polarity components.

Resonance is defined under the constraint:

\[
\frac{\partial E}{\partial P_i} = 0
\]

This formulation is presented as a modeling abstraction rather than a confirmed physiological mechanism.

\section{Interpretive Boundaries}

The polarity field representation introduced here should be interpreted as a formal abstraction. It does not imply established bioelectromagnetic equivalence.

Empirical validation through controlled experimentation remains an essential next step.

\section{Discussion}

An auxiliary experimental direction is outlined in Appendix A, where polarity discrimination is tested using raw optical representations of magnetized samples under double-blind conditions. This protocol is designed to constrain explanatory models by separating direct field exposure from information-mediated conditions.

If classification consistency persists under strictly unprocessed optical data, the nasion-centered polarity model gains structural robustness as an information-sensitive abstraction.

\appendix

\section*{Appendix A: Stage-1 Optical Polarity Discrimination Protocol}

\subsection*{Objective}

To evaluate whether polarity discrimination persists when direct exposure is replaced by photographic representations.

\subsection*{Materials}

\begin{itemize}
\item Magnetized samples (recommended $N \ge 20$)
\item Independently verified polarity orientation
\item Standardized photographic setup
\end{itemize}

\textbf{Raw Data Requirement:}  
All images must be original raw files. No AI enhancement, filtering, contrast adjustment, or post-processing is permitted. Metadata must be preserved.

\subsection*{Statistical Criterion}

For $N=20$:

\[
P(X \ge 15) < 0.05
\]

One-sided binomial test ($\alpha = 0.05$).

For true accuracy $\ge 0.80$, statistical power exceeds 80\%.

\bibliographystyle{plain}
\bibliography{BIMOD_005_references}

\end{document}